\section*{Capítulo 7: La primera imagen de un agujero negro}
\begin{description}[style=multiline,leftmargin=2em,labelwidth=1.5em]

% DOCX paragraph 2071
\item[1.]\hypertarget{note:07-01}{}El plasma\index{plasma|textbf} es uno de los cuatro estados de la materia, junto con el sólido, el líquido y el gaseoso. Se caracteriza por ser una mezcla de iones\index{iones}, es decir, átomos que han perdido algunos o todos sus electrones, y electrones libres. El estado de plasma es muy común en astrofísica debido a las altísimas energías que pueden alcanzarse en el cosmos, por ejemplo, en la superficie de una estrella.

% DOCX paragraph 2072
\item[2.]\hypertarget{note:07-02}{}H. Falcke, M. Kramer, L. Rezzolla, ERC Synergy Grant \emph{BlackHoleCam: Imaging the Event Horizon of Black Holes}; proyecto n.º 610058 (2013-2021).

% DOCX paragraph 2073
\item[3.]\hypertarget{note:07-03}{}El parámetro de impacto se utiliza en física de partículas y en física gravitatoria para medir la distancia perpendicular entre un blanco y el punto más cercano a él que alcanza una partícula a lo largo de su trayectoria.

% DOCX paragraph 2074
\item[4.]\hypertarget{note:07-04}{}Como es natural, dada la simetría esférica\index{simetría esférica} de un agujero negro de Schwarzschild\index{agujero negro!Schwarzschild}, la órbita circular inestable de los fotones es toda una superficie esférica cuyo radio es el del anillo de luz. Este concepto puede extenderse también a los agujeros negros en rotación, pero entonces la situación se complica mucho. Además de distinguir entre trayectorias de fotones en \emph{corrotación} y en \emph{contrarrotación} respecto al agujero negro, hay que sustituir la superficie esférica del anillo de luz por una región que solo es aproximadamente esférica.

% DOCX paragraph 2075
\item[5.]\hypertarget{note:07-05}{}R. Penrose y R. M. Floyd, en \emph{Nature Physical Science}, 229, 177 (1971).

% DOCX paragraph 2076
\item[6.]\hypertarget{note:07-06}{}El efecto que acabamos de describir se denomina \emph{efecto Doppler\index{relativista!efecto Doppler} (relativista) longitudinal}, para distinguirlo del \emph{efecto Doppler relativista transversal}, que produce una variación de la frecuencia incluso cuando el movimiento es perpendicular a la dirección de emisión de la luz. Sin embargo, el efecto Doppler transversal no desempeña un papel importante en lo que estamos explicando aquí.

% DOCX paragraph 2077
\item[7.]\hypertarget{note:07-07}{}En matemáticas, todas estas propiedades se expresan mediante el \emph{grado de no linealidad} de las ecuaciones, que indica cómo una pequeña variación en una propiedad del plasma puede producir enormes variaciones en su comportamiento.

% DOCX paragraph 2078
\item[8.]\hypertarget{note:07-08}{}El positrón es la antipartícula del electrón y, por tanto, tiene la misma masa en reposo que este, pero carga positiva.

% DOCX paragraph 2079
\item[9.]\hypertarget{note:07-09}{}Hay varios mecanismos que podrían proporcionar la energía necesaria para la aceleración. Una posibilidad es que proceda directamente de la energía de rotación del agujero negro, mediante el llamado \emph{proceso de Blandford-Znajek} (véase R. D. Blandford y R. L. Znajek, en \emph{Monthly Notices of the Royal Astronomical Society}, 179, 433 (1977)). Otra es que la proporcione la energía de rotación del disco, a través del \emph{proceso de Blandford-Payne} (véase R. D. Blandford y D. G. Payne, en \emph{Monthly Notices of the Royal Astronomical Society}, 199, 883 (1982)). Por último, la aceleración podría deberse a procesos disipativos dentro del plasma magnetizado.

% DOCX paragraph 2080
\item[10.]\hypertarget{note:07-10}{}T. Karras, S. Laine, T. Aila, \emph{A Style-Based Generator Architecture for Generative Adversarial Networks}, arXiv: 1812.04948.

% DOCX paragraph 2081
\item[11.]\hypertarget{note:07-11}{}Y. Mizuno, Z. Younsi, CM Fromm, O. Porth, M. De Laurentis, H. Olivares, H. Falcke, M. Kramer, L. Rezzolla, en «Nature Astronomy», 2, 585 (2018); H. Olivares, Z. Younsi, C. M. Fromm, M. De Laurentis, O. Porth, Y. Mizuno, H. Falcke, M. Kramer, L. Rezzolla, en \emph{Monthly Notices of the Royal Astronomical Society}, 497, 521 (2020)

\end{description}
